\chapter{El Incidente}
\label{chap:incidente}

\section{Descripción del Atacante}

En este trabajo no se utiliza un adversario existente ni se replica el comportamiento de un APT conocido. En su lugar, se crea un atacante propio con el fin de construir un escenario didáctico totalmente adaptado a los objetivos del ejercicio. Esto permite ajustar, durante la elaboración del ataque, sus tácticas, técnicas y decisiones operativas para que encajen exactamente con las necesidades formativas del proyecto.

Aunque el adversario no está definido por completo al inicio, el ataque presentado sí constituye un escenario cerrado: la intrusión ya ha ocurrido y todas las evidencias están disponibles para su análisis.

\section{Descripción del Ataque (Resumen Ejecutivo)}

El escenario planteado reproduce un ataque estructurado en dos fases, diseñado para simular el comportamiento de un adversario realista. En la fase inicial, el atacante obtiene acceso al equipo víctima mediante un dispositivo físico que le permite ejecutar código de forma inmediata. Con este punto de apoyo establece una conexión remota hacia su infraestructura externa, desde la cual realiza un reconocimiento básico del sistema y prepara la siguiente fase.

En la segunda fase, el adversario despliega su componente principal, diseñado para operar de forma discreta y mantenerse en el sistema sin levantar alertas. Este componente establece persistencia, amplía la capacidad de control del atacante y se integra en procesos legítimos del sistema para dificultar su detección. A partir de este momento, el adversario puede llevar a cabo acciones más avanzadas, como analizar el entorno, escalar privilegios, desplazarse por la red o extraer información.

\section{Mapeo del Ataque a la Cyber Kill Chain y MITRE ATT\&CK}

En esta sección se presenta una referencia técnica que describe el ataque desde el punto de vista del adversario. Se detalla qué acciones llevó a cabo el atacante, cuáles eran sus objetivos en cada fase y qué técnicas MITRE ATT\&CK aplicó. Cada acción se mapea con los artefactos del sistema víctima donde podría encontrarse información relevante sobre dicha actividad, proporcionando una visión clara y útil para el análisis.

La \textit{Cyber Kill Chain} se utiliza como marco para organizar la exposición. El orden elegido es cronológico inverso —del estado final del ataque hacia sus fases iniciales—, ya que este enfoque facilita que el analista comprenda la secuencia real de los hechos: en un escenario DFIR, la investigación siempre comienza desde el impacto final y progresa hacia atrás.

El objetivo de esta sección es ofrecer una guía técnica que relacione de forma directa:
\begin{itemize}
    \item Las acciones ofensivas del atacante,
    \item Las técnicas MITRE ATT\&CK asociadas,
    \item Los artefactos del sistema donde podría hallarse evidencia de cada acción.
\end{itemize}

No se trata de una investigación completa, sino de una referencia estructurada que ayuda a entender cómo se manifiestan las técnicas del atacante en el sistema comprometido.

\subsection{Acciones sobre el objetivo}

\paragraph{Ransomware} En este punto, el atacante ya ha permanecido un periodo considerable dentro de la organización, llevando a cabo diversas actividades de post‑explotación. En la fase final de su operación ejecuta el cifrado de información crítica en el sistema comprometido.

Para una descripción detallada del procedimiento llevado a cabo por el adversario, consúltese el Capítulo~\ref{chap:red-team}, específicamente la Sección~\ref{sec:ejecutar-desde-implante}, donde se analiza el proceso completo con mayor profundidad.
\begin{leftbar}
    \begin{description}
        \item[Táctica:] Impact (TA0040).
        \item[Técnicas:] Data Encrypted for Impact (T1486).
        \item[Procedimiento:] Se desplegó un \textit{.NET assembly} dentro del proceso donde residía el implante. Este componente cifró el fichero \texttt{documento\_importante.txt} utilizando criptografía simétrica. La clave simétrica fue posteriormente cifrada con una clave pública de un esquema asimétrico controlado por el atacante. Finalmente, se dejó una nota de rescate con un correo electrónico de contacto.
    \end{description}
\end{leftbar}

\textbf{Artefactos relevantes}
\begin{itemize}
    \item \textbf{Sysmon Logs}: Buscar carga de assemblies .NET, creación de procesos anómalos, modificaciones de ficheros y eventos de cifrado o acceso masivo.
    \item \textbf{Suricata Logs}: Identificar comunicaciones salientes sospechosas, especialmente relacionadas con la exfiltración de claves o contacto con infraestructura del atacante.
    \item \textbf{Event Logs (Security, Application)}: Revisar ejecución de binarios, eventos de autenticación, fallos o excepciones de aplicaciones .NET y posibles escaladas.
    \item \textbf{Memory Dump}: Localizar el assembly .NET inyectado, claves en memoria, estructuras del implante y cualquier rastro del proceso de cifrado.
    \item \textbf{$MFT$}: Ver cambios en documento\_importante.txt, creación de la nota de rescate y modificaciones de timestamps asociadas al cifrado.
    \item \textbf{$USN\_Journal$}: Identificar operaciones de escritura durante el proceso de cifrado y confirmación de actividad en los ficheros afectados.
    \item \textbf{Prefetch Files}: Determinar ejecución reciente del proceso contenedor del implante y del assembly utilizado.
    \item \textbf{Amcache.hve}: Ver binarios recientes que podrían haber actuado como loader del assembly .NET o como componente auxiliar.
    \item \textbf{ShimCache}: Confirmar ejecución histórica del binario que alojaba el implante o de herramientas usadas por el atacante.
    \item \textbf{.NET CLR Logs}: Buscar carga dinámica de assemblies, errores, o eventos de ejecución sospechosos relacionados con el cifrado.
    \item \textbf{Antivirus/EDR Logs}: Detectar alertas, bloqueos parciales o telemetría sobre inyección de código, cifrado o movimiento lateral.
    \item \textbf{Ransom Note Metadata}: Analizar timestamps, usuario que la generó y posible correlación con la fase final del ataque.
\end{itemize}

\subsection{Mando y Control (C2)}

\paragraph{Sliver C2} En este punto, el atacante ya ha superado las fases de explotación e instalación y dispone de un mecanismo persistente que le permite interactuar de forma continua con el sistema víctima mediante un implante de \textit{Sliver}.

Para una descripción detallada del proceso mediante el cual el implante de Sliver se carga en memoria en cada ejecución, evitando la detección por parte del EDR, consúltese el Capítulo~\ref{chap:red-team}, Sección~\ref{sec:reflective}. 

Asimismo, para obtener información sobre la infraestructura de servidores utilizada por el atacante, véase la Sección~\ref{sec:infraestructura}.

\begin{leftbar}
    \begin{description}
        \item[Táctica:] Defense Evasion (TA0005).
        \item[Técnica:] Reflective Code Loading (T1620).
        \item[Procedimiento:] La DLL maliciosa utiliza carga reflexiva para ejecutarse dentro del proceso comprometido. El módulo \texttt{calibre-launcher.dll} establece una conexión con el servidor C2, descarga el implante de Sliver y lo carga en memoria mediante un mecanismo de reflexión, sin utilizar las APIs estándar de carga de librerías ni generar artefactos adicionales en disco.
    \end{description}
\end{leftbar}

\textbf{Artefactos relevantes}
\begin{itemize}
    \item \textbf{Sysmon Logs}: Buscar eventos de conexión saliente (Event ID 3), creación de hilos remotos, acceso a memoria de otros procesos y carga de módulos no estándar (Event ID 7 con anomalías en la ruta).
    \item \textbf{Suricata Logs}: Identificar patrones de C2, tráfico persistente o no habitual asociado al framework Sliver y a la fase de descarga del implante.
    \item \textbf{Event Logs (Security, Application)}: Revisar eventos de ejecución irregular, fallos de módulos .NET o errores de carga que puedan correlacionarse con la inyección reflexiva.
    \item \textbf{Memory Dump}: Localizar el implante Sliver cargado en memoria, regiones marcadas como ejecutables sin mapeo en disco y código asociado a la carga reflexiva.
    \item \textbf{$MFT$}: Verificar la ausencia de escritura de DLLs o ficheros relacionados para confirmar ejecución fileless.
    \item \textbf{$USN\_Journal$}: Corroborar que no se han generado artefactos en disco durante la infección (indicador de carga reflexiva).
    \item \textbf{Prefetch Files}: Identificar el proceso anfitrión utilizado para la carga reflexiva (solo evidencia de ejecución del proceso vulnerable).
    \item \textbf{Amcache.hve}: Confirmar si calibre-launcher.dll o componentes auxiliares llegaron a ejecutarse desde disco antes de la fase fileless.
    \item \textbf{ShimCache}: Ver cualquier huella de ejecución histórica del proceso legítimo comprometido.
    \item \textbf{Antivirus/EDR Logs}: Revisar detecciones de inyección, anomalías de memoria, creación de secciones ejecutables o comportamiento asociado a Sliver.
    \item \textbf{Network Connections Artifacts}: Buscar conexiones persistentes o repetitivas hacia el C2 del atacante.
\end{itemize}

\begin{leftbar}
    \begin{description}
        \item[Táctica:] Command and Control (TA0011).
        \item[Técnicas:]
        \begin{itemize} 
            \item Application Layer Protocol: Web Protocols (T1071.001)
            \item Encrypted Channel (T1573)
        \end{itemize}
        \item[Procedimiento:]  
        El implante residente en memoria mantiene un canal de comunicación persistente con la infraestructura de mando y control mediante un mecanismo de \textit{beaconing}. Para reducir la probabilidad de detección, el tráfico se encapsula dentro de peticiones HTTPS que imitan patrones legítimos de navegación web. Además, todas las comunicaciones viajan cifradas, lo que impide la inspección de su contenido y permite al atacante intercambiar órdenes y resultados sin generar artefactos visibles o anomalías en el tráfico.
    \end{description}
\end{leftbar}

\textbf{Artefactos relevantes}
\begin{itemize}
    \item \textbf{Sysmon Logs}: Buscar conexiones salientes periódicas (Event ID 3), intervalos regulares de beaconing y patrones anómalos de uso de HTTPS desde procesos que no deberían generar dicho tráfico.
    \item \textbf{Suricata Logs}: Detectar dominios o IPs sospechosas, JA3/JA3S que no coinciden con navegadores legítimos y patrones repetitivos de petición-respuesta propios de beaconing.
    \item \textbf{Event Logs (Security, Application)}: Revisar fallos o eventos de red inusuales en procesos que no suelen comunicarse externamente.
    \item \textbf{Memory Dump}: Identificar el implante en memoria, estructuras de comunicación cifrada, temporizadores de beaconing y buffers utilizados para HTTPS.
    \item \textbf{Network Connections Artifacts}: Verificar conexiones persistentes hacia los mismos endpoints y correlacionar tiempos de comunicación.
    \item \textbf{Antivirus/EDR Logs}: Buscar telemetría sobre conexiones repetitivas, anomalías de comportamiento en procesos y posibles indicadores de tráfico camuflado.
\end{itemize}

\subsection{Instalación}

\paragraph{Stage 2} En esta fase, el atacante ya ha superado la etapa de explotación inicial y ha logrado ejecutar código en el sistema víctima, obteniendo una \textit{reverse shell}. A partir de este acceso, procede a preparar el \textit{stage 2} del ataque y a desplegar su \textit{payload} final, cuyo objetivo es establecer un mecanismo de acceso persistente al sistema comprometido.

Para lograrlo, el adversario emplea una combinación de técnicas basadas en \textit{DLL Hijacking}, \textit{Side Loading} y \textit{DLL Proxying}. Un análisis detallado de este procedimiento puede consultarse en el Capítulo~\ref{chap:red-team}, Sección~\ref{sec:dll-hijacking}.

\begin{leftbar}
    \begin{description}
        \item[Táctica:] Defense Evasion (TA0005).
        \item[Técnica:] Modify Registry or File Permissions: File and Directory Permissions Modification (T1222).
        \item[Procedimiento:]  
        El atacante copia el ejecutable legítimo de Calibre desde su ubicación original a un directorio donde dispone de permisos de escritura, concretamente \texttt{\%LOCALAPPDATA\%\textbackslash Microsoft\textbackslash WindowsApps}. Esta acción le permite controlar el entorno de carga de librerías del binario y establecer las condiciones necesarias para llevar a cabo el secuestro de DLL.
    \end{description}
\end{leftbar}

\textbf{Artefactos relevantes}
\begin{itemize}
    \item \textbf{$MFT$}: Identificar la copia del binario legítimo de Calibre y los cambios de ubicación asociados al secuestro de DLL.
    \item \textbf{$USN\_Journal$}: Corroborar operaciones de copia, creación y modificación dentro de \%LOCALAPPDATA\%\\Microsoft\\WindowsApps.
    \item \textbf{Sysmon Logs}: Revisar eventos de creación de ficheros (Event ID 11), carga de módulos sospechosos (Event ID 7) y ejecución del binario adulterado.
    \item \textbf{Suricata Logs}: Ver si, tras el secuestro, el binario modificado genera tráfico anómalo o relacionado con C2.
    \item \textbf{Prefetch Files}: Verificar ejecución del binario de Calibre desde la nueva ruta controlada por el atacante.
    \item \textbf{Amcache.hve}: Confirmar información sobre el binario copiado, su ruta y su primera ejecución desde el directorio comprometido.
    \item \textbf{ShimCache}: Detectar huellas de ejecución del binario modificado, útil para correlacionar la activación del DLL hijacking.
    \item \textbf{Event Logs (Security)}: Revisar eventos de acceso a ficheros, permisos y posibles modificaciones en directorios donde el usuario no debería trabajar normalmente.
\end{itemize}

\begin{leftbar}
    \begin{description}
        \item[Táctica:] Execution (TA0002).
        \item[Técnica:] Ingress Tool Transfer (T1105).
        \item[Procedimiento:]  
        El atacante descarga en el directorio del binario legítimo una DLL maliciosa denominada \texttt{calibre-launcher.dll}, diseñada específicamente para reemplazar a la librería legítima. Esta DLL ya contiene la carga maliciosa integrada, lista para ejecutar el implante en cuanto la aplicación intente cargarla de forma normal.
    \end{description}
\end{leftbar}

\textbf{Artefactos relevantes}
\begin{itemize}
    \item \textbf{$MFT$}: Identificar la creación de calibre-launcher.dll maliciosa en el directorio del binario.
    \item \textbf{$USN\_Journal$}: Ver operaciones de escritura y sustitución de la DLL legítima por la maliciosa.
    \item \textbf{Sysmon Logs}: Revisar eventos de creación de ficheros (Event ID 11) y carga de la DLL manipulada (Event ID 7) por el ejecutable de Calibre.
    \item \textbf{Amcache.hve}: Confirmar registro de la DLL maliciosa si llegó a ejecutarse o cargarse desde disco.
    \item \textbf{ShimCache}: Ver huellas de ejecución del binario que carga la DLL maliciosa.
    \item \textbf{Prefetch Files}: Validar que el ejecutable legítimo fue ejecutado desde el directorio alterado, provocando la carga de la DLL maliciosa.
    \item \textbf{Event Logs (Security)}: Revisar eventos relacionados con acceso, modificación o escritura en el directorio del binario.
    \item \textbf{Antivirus/EDR Logs}: Buscar alertas sobre carga de DLL no firmada, inyección o comportamiento sospechoso.
\end{itemize}

\begin{leftbar}
    \begin{description}
        \item[Táctica:] Defense Evasion / Privilege Escalation (TA0005 / TA0004).
        \item[Técnica:] 
        \begin{itemize} 
            \item Hijack Execution Flow: DLL Search Order Hijacking (T1574.001).
            \item Masquerading: Match Legitimate Name or Location (T1036.005).
        \end{itemize} 
        \item[Procedimiento:]  
        El atacante renombra la DLL legítima a \texttt{calibre-launcher-old.dll} para conservar su funcionalidad original. A continuación, coloca en el directorio una DLL maliciosa con el nombre legítimo esperado por la aplicación (\texttt{calibre-launcher.dll}). Debido al orden de búsqueda de DLLs en Windows, la aplicación carga primero la DLL maliciosa, que actúa como un proxy reenviando las llamadas a la DLL legítima renombrada mientras ejecuta su código malicioso.
    \end{description}
\end{leftbar}

\textbf{Artefactos relevantes}
\begin{itemize}
    \item \textbf{$MFT$}: Identificar renombrado de la DLL legítima y creación de la DLL maliciosa con el nombre esperado.
    \item \textbf{$USN\_Journal$}: Corroborar operaciones de renombrado, escritura y sustitución dentro del directorio.
    \item \textbf{Sysmon Logs}: Revisar creación y modificación de ficheros (Event ID 11) y carga de calibre-launcher.dll maliciosa (Event ID 7).
    \item \textbf{Prefetch Files}: Validar ejecución del proceso que provoca la carga de la DLL maliciosa como primera en el orden de búsqueda.
    \item \textbf{Amcache.hve}: Confirmar la presencia y ejecución previa de calibre-launcher.dll y calibre-launcher-old.dll.
    \item \textbf{ShimCache}: Detectar huellas de carga tanto de la DLL maliciosa como de la DLL legítima renombrada.
    \item \textbf{Event Logs (Security)}: Ver eventos de modificación y escritura sobre el directorio del binario.
    \item \textbf{Antivirus/EDR Logs}: Buscar alertas por DLL hijacking, anomalías en la cadena de carga o uso de una DLL acting-as-proxy.
\end{itemize}

\paragraph{Persistencia} El atacante establece un mecanismo de persistencia mediante la creación de una tarea programada, garantizando así que su \textit{payload} se ejecute automáticamente tras cada reinicio del sistema y mantenga el acceso al entorno comprometido.

Para una descripción detallada del procedimiento, consúltese el Capítulo~\ref{chap:red-team}, Sección~\ref{sec:persistencia}.

\begin{leftbar}
    \begin{description}
        \item[Táctica:] Persistence (TA0003).
        \item[Técnica:] Scheduled Task/Job (T1053).
        \item[Procedimiento:]  
        El atacante crea una tarea programada denominada \texttt{calibre-update}, configurada para ejecutarse en cada reinicio del sistema. Esta tarea invoca \texttt{calibre-debug.exe}, un ejecutable legítimo presente en la instalación original y que puede ejecutarse sin interfaz gráfica. Al lanzarse automáticamente, el binario vuelve a cargar la DLL maliciosa, permitiendo que el implante de Sliver recupere la persistencia tras cada reinicio.
    \end{description}
\end{leftbar}

\textbf{Artefactos relevantes}
\begin{itemize}
    \item \textbf{Scheduled Tasks (schtasks.xml / TaskCache)}: Identificar la creación de la tarea calibre-update y su configuración de persistencia.
    \item \textbf{Event Logs (Microsoft-Windows-TaskScheduler/Operational)}: Buscar eventos de creación, modificación y ejecución de la tarea programada.
    \item \textbf{Sysmon Logs}: Revisar ejecuciones automáticas de calibre-debug.exe (Event ID 1) y carga de la DLL maliciosa (Event ID 7).
    \item \textbf{Prefetch Files}: Validar que calibre-debug.exe se ejecuta tras cada reinicio.
    \item \textbf{$MFT$}: Confirmar la presencia de calibre-debug.exe y de la DLL maliciosa en el directorio utilizado.
    \item \textbf{$USN\_Journal$}: Ver operaciones de escritura o modificación relacionadas con la creación de la tarea o binarios invocados.
    \item \textbf{Amcache.hve}: Registrar ejecuciones previas de calibre-debug.exe desde la ubicación alterada.
    \item \textbf{ShimCache}: Identificar huellas de ejecución del binario encargado de activar la persistencia.
    \item \textbf{Antivirus/EDR Logs}: Buscar detecciones relacionadas con persistencia, actividad automatizada o carga de DLL maliciosa.
\end{itemize}

\subsection{Explotación}

\paragraph{Stage 1} En este punto, el atacante obtiene la primera ejecución de código en el sistema víctima y despliega el \textit{payload} inicial que le permite establecer una reverseshell con el equipo comprometido.

Para una descripción detallada del desarrollo y funcionamiento de este \textit{payload}, véase el Capítulo~\ref{chap:red-team}, Sección~\ref{sec:shellcode}.

\begin{leftbar}
    \begin{description}
        \item[Táctica:] Execution (TA0002) / Command and Control (TA0011).
        \item[Técnica:]
        \begin{description}
            \item Command and Scripting Interpreter: PowerShell (T1059.001).
            \item Ingress Tool Transfer (T1105).
            \item User Execution: Malicious File (T1204.002).
            \item Remote Services / Reverse Shell (T1021).
        \end{description}
        \item[Procedimiento:]  
        El atacante ejecuta un script de PowerShell que establece conexión con la infraestructura C2 y descarga un archivo temporal denominado \texttt{2fd19d11-d4b7-46d1-94db-7edd16980d15.tmp} en el directorio temporal del usuario (\texttt{\$env:TEMP}).  
        A continuación, renombra el archivo con extensión \texttt{.exe} y lo ejecuta de forma oculta mediante \texttt{-WindowStyle Hidden}. El binario lanzado establece una reverse shell hacia el servidor controlado por el atacante, brindándole acceso remoto al sistema comprometido.
    \end{description}
\end{leftbar}

\textbf{Artefactos relevantes}
\begin{itemize}
    \item \textbf{PowerShell Operational Log}: Identificar ejecución del script, comandos usados, descarga del fichero y parámetros como \texttt{-WindowStyle Hidden}.
    \item \textbf{Sysmon Logs}: Revisar eventos de creación de procesos (Event ID 1), creación de ficheros temporales (Event ID 11) y conexiones salientes asociadas a la reverse shell (Event ID 3).
    \item \textbf{Suricata Logs}: Detectar tráfico hacia el C2, patrones de reverse shell y comunicaciones anómalas asociadas al binario descargado.
    \item \textbf{$MFT$}: Confirmar creación del archivo temporal y su posterior renombrado a extensión .exe.
    \item \textbf{$USN\_Journal$}: Corroborar operaciones de escritura y renombrado dentro de \texttt{\$env:TEMP}.
    \item \textbf{Event Logs (Security)}: Revisar eventos asociados a ejecución de scripts, permisos y posibles actividades del proceso resultante.
    \item \textbf{Prefetch Files}: Validar la ejecución del binario renombrado desde el directorio temporal.
    \item \textbf{Amcache.hve}: Registrar la primera ejecución del binario .exe descargado.
    \item \textbf{ShimCache}: Identificar huellas de ejecución tanto del script de PowerShell como del binario temporal.
    \item \textbf{Memory Dump}: Localizar artefactos del proceso de reverse shell, buffers de comunicación y posibles credenciales en memoria.
    \item \textbf{Antivirus/EDR Logs}: Buscar alertas relacionadas con ejecución de PowerShell, creación de binarios en TEMP y conexiones de remote access.
\end{itemize}

\subsection{Entrega}

\paragraph{BadUSB} En este punto, el atacante interactúa por primera vez con la organización víctima, marcando el inicio del incidente de seguridad. El objetivo de esta acción es introducir el primer \textit{payload} dentro del entorno de la víctima mediante un dispositivo \textit{BadUSB}.

Para un análisis detallado de las decisiones tomadas durante esta fase, véase el Capítulo~\ref{chap:red-team}, Sección~\ref{sec:bad-usb}.

\begin{leftbar}
    \begin{description}
        \item[Táctica:] Initial Access (TA0001) / Execution (TA0002).
        \item[Técnica:] 
        \begin{description}
            \item Replication Through Removable Media (T1091)
            \item Input Injection (T1674)
        \end{description}
        \item[Procedimiento:]  
        El atacante conecta un dispositivo \textit{BadUSB} que emula un teclado (HID). Nada más inicializarse, el dispositivo inyecta pulsaciones de teclado que abren el cuadro de ejecución mediante \texttt{Win + R} y lanzan una consola de PowerShell. Desde ella, se ejecutan los comandos necesarios para descargar y ejecutar el script malicioso sin interacción del usuario.
    \end{description}
\end{leftbar}

\textbf{Artefactos relevantes.}  
\begin{itemize}
    \item \textbf{Event Logs (Security)}: Buscar eventos de conexión de dispositivos USB y actividad asociada a entradas HID simuladas.
    \item \textbf{Sysmon Logs}: Revisar creación de procesos (Event ID 1) y apertura de PowerShell iniciada sin interacción del usuario.
    \item \textbf{PowerShell Operational Log}: Identificar comandos ejecutados automáticamente tras la inyección de pulsaciones.
    \item \textbf{$MFT$}: Confirmar la descarga de cualquier fichero relacionado con el script malicioso.
    \item \textbf{$USN\_Journal$}: Corroborar operaciones de escritura en el directorio donde se descargue el script.
    \item \textbf{Suricata Logs}: Detectar tráfico generado por la descarga o ejecución remota del script.
    \item \textbf{Prefetch Files}: Validar ejecución automática de PowerShell y de cualquier binario descargado.
    \item \textbf{Amcache.hve}: Registrar ejecución del script o binarios asociados.
    \item \textbf{ShimCache}: Identificar huellas de ejecución del proceso PowerShell iniciado por las pulsaciones inyectadas.
    \item \textbf{Antivirus/EDR Logs}: Detectar actividad anómala de PowerShell o ejecución sin interacción del usuario.
\end{itemize}

\subsection{Armamento}

En esta fase, el atacante prepara los \textit{payloads}, herramientas y la infraestructura que utilizará durante el ataque. Al igual que ocurre con la fase de reconocimiento, esta etapa se sitúa más cerca del análisis preventivo que del análisis forense de un incidente.

Las acciones llevadas a cabo por el atacante en esta fase se describen en el capítulo \textit{Red Team Notes} (ver Sección~\ref{chap:red-team}), donde se detallan, desde la perspectiva ofensiva, el desarrollo de los \textit{payloads}, las decisiones técnicas tomadas y la infraestructura empleada para ejecutar el ataque.

Desde el punto de vista del analista, esta fase puede apoyarse en servicios externos que realizan escaneos continuos de Internet en busca de infraestructuras maliciosas, como servidores de \textit{Cobalt Strike} u otros \textit{C2} con configuraciones mínimas o pobre \textit{fingerprinting}. Estos servicios permiten a la organización detectar tempranamente posibles activos maliciosos relacionados con campañas activas y, en caso necesario, aplicar medidas preventivas como el bloqueo automático de conexiones hacia o desde dichas infraestructuras.

\subsection{Reconocimiento}

En esta fase, el atacante intenta recopilar información relevante sobre la organización, sus sistemas y posibles vectores de entrada. Desde el punto de vista del analista, estas actividades suelen considerarse parte de un análisis preventivo más que de un análisis de incidente propiamente dicho.

Dentro de esta categoría se incluyen tareas como el monitoreo de actividades de reconocimiento dirigidas a los endpoints expuestos de la organización, la revisión de bases de datos comercializadas por ciberdelincuentes que contengan información relacionada con la empresa y, en general, la supervisión de cualquier actividad externa que pueda indicar la recopilación de información previa a un ataque.